\begin{abstract}
Vision-language models (VLMs) often treat textual context as a shortcut, even when that text conflicts with visual evidence.
This ``blind faith in text'' is especially problematic in document and visual question answering, where auxiliary text may be incomplete, stale, adversarial, or automatically retrieved.
We study whether preference optimization can reduce this failure mode without sacrificing normal visual question answering performance.
The goal is not to teach the model to ignore all text: auxiliary text is often useful when it agrees with the image, and the model usually does not know in advance whether the text is reliable.
Starting from corrupted-text evaluations, we mine model-specific text-following errors: cases where the model is wrong and its wrong answer is copied from the corrupted auxiliary text.
We then train VLMs with Direct Preference Optimization (DPO), preferring the image-grounded answer over the model's own text-following error.
On DocVQA with Qwen3-VL-8B, text-following span DPO substantially improves corrupted-text robustness and reduces the fraction of remaining errors copied from misleading text.
A multi-condition error-mined extension further strengthens this effect.
By contrast, GRPO-style sampled-answer optimization and auxiliary SFT preservation do not improve over DPO.
Experiments on VQAv2 and transfer to Qwen2-VL and LLaVA-Next show that the method works best when the base VLM already has sufficient visual and OCR competence.
These results suggest that high-precision preference pairs are an effective and practical mechanism for correcting text bias in VLMs, while also revealing clear boundary conditions for the approach.
\end{abstract}
